\subsection{Lax localizations}
\label{subsec:laxloc-def}

\gutentag{00C8}%
An \emph{adjunction of $(\infty,2)$-categories} is an adjunction
(\cref{def:adjunction}) in $\twoCat$, regarded as an $(\infty,2)$-category with
hom-categories
$\iota_{1}\func(\A,\B)$:
the non-invertible $2$-morphisms of $\func(\A,\B)$ are discarded, not inverted.

\gutentag{00C9}%
Throughout the rest of this section, $L \dashv R$ is an adjunction of
$(\infty,2)$-categories, with
$L\colon \D \to \C$,
$R\colon \C \to \D$
and unit
$\eta\colon \mathrm{id}_{\D} \Rightarrow RL$.
We write
$\varepsilon\colon LR \Rightarrow \mathrm{id}_{\C}$
for its counit and
\begin{equation*}
\varphi\colon \varepsilon L \circ L\eta \simeq \mathrm{id}_{L} \ , \qquad
\theta\colon R\varepsilon \circ \eta R \simeq \mathrm{id}_{R}
\end{equation*}
for the invertible $2$-morphisms of $\func(\D,\C)$ and $\func(\C,\D)$ given by
the triangle identities.
For $d \in \D$ and $c \in \C$ the composite
\begin{equation}
\label{eq:laxloc-adjunction}
\Hom_{\C}(Ld,c) \xrightarrow{\ R\ } \Hom_{\D}(RLd,Rc)
\xrightarrow{\ \eta_{d}^{\natural}\ } \Hom_{\D}(d,Rc)
\end{equation}
is an equivalence, natural in $d$ and $c$ as functors
$\D\opcat \times \C \to \Cat$:
an inverse is
$g \mapsto \varepsilon_{c} \circ Lg$,
by naturality of $\eta$ and
$\varepsilon$ and by $\varphi$ and $\theta$.
Restricting the adjunction
$\Adj \to \twoCat$
to the full sub-$(\infty,2)$-category of $\Adj$ on $-$, the walking monad
\cite[Not.~5.1]{haugseng2021laxmonads}, gives a monad $T := RL$ on $\D$
(\cref{def:monad}) by \cite[Cor.~8.9]{haugseng2021laxmonads}, with unit
$\eta$, multiplication
$\mu := R\varepsilon L$
and unit laws
\begin{equation*}
\theta L\colon \mu \circ \eta T \simeq \mathrm{id}_{T} \ , \qquad
R\varphi\colon \mu \circ T\eta \simeq \mathrm{id}_{T} \ .
\end{equation*}

\begin{defi}
\label{def:laxloc}
\gutentag{00CA}%
The adjunction $L \dashv R$ is a \emph{lax localization} if $T$ is lax
idempotent, and a \emph{colax localization} if $T$ is colax idempotent
(\cref{def:kz}).
\end{defi}

\begin{lem}
\label{lem:laxloc-criterion}
\gutentag{00CB}%
Let $f\colon x \to y$ and $u\colon y \to x$ be $1$-morphisms of an
$(\infty,2)$-category $\A$.
\begin{enumerate}
\item Given an invertible
$c\colon fu \simeq \mathrm{id}_{y}$,
we have $f \dashv u$ with counit $c$ if and only if there is
$\nu\colon \mathrm{id}_{x} \Rightarrow uf$
with $f\nu$ and $\nu u$ invertible.
\item Given an invertible
$\nu\colon \mathrm{id}_{x} \simeq uf$,
we have $f \dashv u$ with unit $\nu$ if and only if there is
$c\colon fu \Rightarrow \mathrm{id}_{y}$
with $cf$ and $uc$ invertible.
\end{enumerate}
\end{lem}

\begin{proof}
\gutentag{00CC}%
(1) The unit of an adjunction $f \dashv u$ with counit $c$ is such a $\nu$,
as $c$ is invertible, by
the triangle identities. Conversely, by \cite[Cor.~2.2]{bourke2025recognition},
a consequence of \cite[Lemma~2.1.11]{riehlverity2022elements}, a $\nu$ as in
(1) yields a $2$-morphism
$\mathrm{id}_{x} \Rightarrow uf$
satisfying the triangle identities with $c$.
\todo{2-categorical input: check that \cite[Cor.~2.2]{bourke2025recognition} works verbatim.} By \cref{thm:adj-lifting}(2), these extend to an adjunction
$f \dashv u$ with counit $c$. (2) is (1) in
$\A^{\mathrm{co}}$.
\end{proof}

\begin{lem}
\label{lem:laxloc-equiv}
\gutentag{00CD}%
The following are equivalent:
\begin{enumerate}
\item $L \dashv R$ is a lax localization;
\item
$L\eta \dashv \varepsilon L$
in $\func(\D,\C)$, with unit $\varphi^{-1}$;
\item
$R\varepsilon \dashv \eta R$
in $\func(\C,\D)$, with counit $\theta$;
\item
$L\eta_{d} \dashv \varepsilon_{Ld}$
with unit $\varphi_{d}^{-1}$, for
every $d \in \D$;
\item
$R\varepsilon_{c} \dashv \eta_{Rc}$
with counit $\theta_{c}$, for every
$c \in \C$.
\end{enumerate}
In particular, for a lax localization, $L\eta$ is coreflective in
$\func(\D,\C)$ and $\eta R$ is reflective in $\func(\C,\D)$
(\cref{def:reflection}), and so are their components $L\eta_{d}$ and
$\eta_{Rc}$.
\end{lem}

\begin{proof}
\gutentag{00CE}%
Evaluation preserves adjunctions, so (2) implies (4) and (3) implies (5).
Postcomposition with $R$ takes (2) to
$T\eta \dashv \mu$
with unit $(R\varphi)^{-1}$, and precomposition with $L$ takes (3) to
$\mu \dashv \eta T$
with counit $\theta L$; so (2) and (3) each imply (1), by \cref{def:kz} and
\cref{lem:kz-basic}(1).

Assume (1). Let $\upsilon$ be the unit of
$\mu \dashv \eta T$
(\cref{lem:kz-basic}(1)) and
$\delta\colon T\eta \Rightarrow \eta T$
the $2$-morphism $\upsilon T\eta$ followed by the unit law, as in
\cref{prop:kz-action}. By the triangle identities, as the counit of
$\mu \dashv \eta T$
is invertible, so are $\mu\upsilon$ and $\upsilon\eta T$. Hence $\mu\delta$ is
invertible, and so is $\delta_{d}\eta_{d}$: up to invertible $2$-morphisms, it
is $\upsilon_{d}$ whiskered with
$T\eta_{d} \circ \eta_{d} \simeq \eta_{Td} \circ \eta_{d}$.

(1)$\Rightarrow$(5). Let $x := Rc$ and
$a := R\varepsilon_{c}$,
so that
$\theta_{c}\colon a\eta_{x} \simeq \mathrm{id}_{x}$,
and let $\nu$ be the composite
\begin{equation*}
\mathrm{id}_{Tx} \simeq Ta \circ T\eta_{x} \xRightarrow{\ Ta\,\delta_{x}\ } Ta
\circ \eta_{Tx} \simeq \eta_{x} \circ a \ .
\end{equation*}
Then $\nu\eta_{x}$ is invertible as $\delta_{x}\eta_{x}$ is, and $a\nu$ is
invertible as $a\mu_{x}\delta_{x}$ is, since
$a \circ Ta \simeq a \circ \mu_{x}$
by $R$ applied to the naturality of $\varepsilon$ at $\varepsilon_{c}$.
So $a \dashv \eta_{x}$ with counit $\theta_{c}$ by
\cref{lem:laxloc-criterion}(1).

(1)$\Rightarrow$(4). The equivalence \eqref{eq:laxloc-adjunction}
$\Hom_{\C}(LTd,LTd) \simeq \Hom_{\D}(Td,TTd)$
sends
$L\eta_{d} \circ \varepsilon_{Ld}$
to
$T\eta_{d} \circ \mu_{d} \circ \eta_{Td} \simeq T\eta_{d}$
and $\mathrm{id}_{LTd}$ to $\eta_{Td}$, so $\delta_{d}$ corresponds to a
$2$-morphism
$\kappa\colon L\eta_{d} \circ \varepsilon_{Ld} \Rightarrow \mathrm{id}_{LTd}$.
By naturality of \eqref{eq:laxloc-adjunction},
$\varepsilon_{Ld}\kappa$
corresponds to $\mu_{d}\delta_{d}$ and $\kappa L\eta_{d}$ to
$\delta_{d}\eta_{d}$, so both are invertible, and
$L\eta_{d} \dashv \varepsilon_{Ld}$
with unit $\varphi_{d}^{-1}$ by \cref{lem:laxloc-criterion}(2).

(5)$\Rightarrow$(3). By
\cite[Cor.~5.2.12]{abellangagnahaugseng2024straightening},
$R\varepsilon$ is a left adjoint in $\func(\C,\D)$ if its components are and
its naturality squares are right Beck--Chevalley (\cref{lem:mates}(1)).
At $g\colon c \to c'$ the right mate is
\begin{equation*}
TRg \circ \eta_{Rc} \Longrightarrow \eta_{Rc'}R\varepsilon_{c'} \circ TRg \circ
\eta_{Rc} \simeq \eta_{Rc'} \circ Rg \circ R\varepsilon_{c} \circ \eta_{Rc}
\overset{\theta_{c}}{\simeq} \eta_{Rc'} \circ Rg \ ,
\end{equation*}
whose first $2$-morphism is the unit of
$R\varepsilon_{c'} \dashv \eta_{Rc'}$
at
$TRg \circ \eta_{Rc} \simeq \eta_{Rc'} \circ Rg$,
invertible by a triangle identity as the counit is invertible. So
$R\varepsilon \dashv G$
for some $G$. The $2$-morphism
$\eta R \Rightarrow G$
adjunct to $\theta$ has as components the comparisons between the right adjoints
$\eta_{Rc}$ and $G_{c}$ of $R\varepsilon_{c}$, so it is invertible, as
$2$-morphisms of $\func(\C,\D)$ are detected componentwise (see the proof of
\cref{lem:mates-refl}). Hence
$R\varepsilon \dashv \eta R$
with counit $\theta$.

(4)$\Rightarrow$(2). Likewise, at $g\colon d \to d'$ the right mate of the
naturality square of $L\eta$ is
\begin{equation*}
Lg \circ \varepsilon_{Ld} \xRightarrow{\ \varphi_{d'}^{-1}\ }
\varepsilon_{Ld'}L\eta_{d'} \circ Lg \circ \varepsilon_{Ld} \simeq
\varepsilon_{Ld'} \circ LTg \circ L\eta_{d} \circ \varepsilon_{Ld}
\Longrightarrow \varepsilon_{Ld'} \circ LTg \ ,
\end{equation*}
whose last $2$-morphism is the counit $\kappa_{d}$ of
$L\eta_{d} \dashv \varepsilon_{Ld}$
whiskered with
$\varepsilon_{Ld'} \circ LTg \simeq Lg \circ \varepsilon_{Ld}$;
it is invertible as
$\varepsilon_{Ld}\kappa_{d}$
is, by a triangle identity.
So $L\eta \dashv G'$ for some $G'$, and the $2$-morphism
$G' \Rightarrow \varepsilon L$
adjunct to $\varphi^{-1}$ is invertible, as its components are the comparisons
between the right adjoints $G'_{d}$ and $\varepsilon_{Ld}$ of $L\eta_{d}$.
\end{proof}

\begin{rmk}
\label{rmk:laxloc-literature}
\gutentag{00CF}%
For $2$-categories, Bourke \cite[\S 2.1]{bourke2025recognition} calls a
biadjunction lax-idempotent if (3) holds, and the equivalence of (1) and (3) is
\cite[Prop.~2.5.12]{stepan2025thesis}. For $(\infty,2)$-categories we know of no
reference (\cref{rmk:kz-literature}).
\end{rmk}

\gutentag{00D0}%
As $(-)^{\mathrm{co}}$ is an automorphism of $\twoCat$ with
$\iota_{1}\func(\A^{\mathrm{co}},\B^{\mathrm{co}}) \simeq
\iota_{1}\func(\A,\B)$
compatibly with composition (\cref{subsec:prelim-monads}), an adjunction
$L \dashv R$ is also an adjunction
$L^{\mathrm{co}} \dashv R^{\mathrm{co}}$
between $\D^{\mathrm{co}}$ and $\C^{\mathrm{co}}$, with the same unit, counit
and monad. By \cref{def:kz}, $L \dashv R$ is a colax localization if
and only if
$L^{\mathrm{co}} \dashv R^{\mathrm{co}}$
is a lax localization. We treat lax localizations, and will spell out the colax
case in \cref{rmk:laxloc-colax}.
